\chapter{วิวัฒนาการและความสำคัญของจุลชีววิทยาทางการเกษตร}
\label{chap:intro}

% ==========================================================
% แผนบริหารการสอนประจำบท (Chapter Syllabus)
% ==========================================================
\section*{แผนบริหารการสอนประจำบทนำ}
\addcontentsline{toc}{section}{แผนบริหารการสอนประจำบทนำ}

\textbf{หัวข้อเนื้อหาประจำบท}
\begin{enumerate}
    \item ความหมายและขอบข่ายของจุลชีววิทยาสำหรับการเกษตร
    \item ลำดับเวลาและหมุดหมายสำคัญทางประวัติศาสตร์ของจุลชีววิทยาการเกษตร
    \item ความเชื่อมโยงระหว่างจุลินทรีย์กับความมั่นคงทางอาหารและสิ่งแวดล้อม
    \item บทบาทของจุลินทรีย์ในการขับเคลื่อนโมเดลเศรษฐกิจชีวภาพ (BCG Model)
\end{enumerate}

\textbf{วัตถุประสงค์เชิงพฤติกรรม}
\begin{itemize}
    \item อธิบายขอบข่ายและความสำคัญของจุลชีววิทยาในระบบนิเวศการเกษตรได้
    \item ระบุและวิเคราะห์ผลงานของนักวิทยาศาสตร์ผู้บุกเบิกในอดีตที่ส่งผลต่อการเกษตรสมัยใหม่ได้
    \item อธิบายบทบาทของจุลินทรีย์ในการสนับสนุนเป้าหมายการพัฒนาที่ยั่งยืนและเศรษฐกิจชีวภาพได้
\end{itemize}

\textbf{กิจกรรมการเรียนการสอน}
\begin{itemize}
    \item การบรรยายเชิงวิชาการร่วมกับการอภิปรายกรณีศึกษาพัฒนาการทางประวัติศาสตร์
    \item การศึกษาและสรุปแผนผังเส้นเวลาของนวัตกรรมชีวภาพทางการเกษตร
    \item การระดมความคิดเห็นเกี่ยวกับแนวทางการนำจุลินทรีย์มาทดแทนสารเคมีในท้องถิ่นภาคตะวันออก
\end{itemize}

\textbf{สื่อการเรียนการสอน}
\begin{itemize}
    \item เอกสารคำสอนและภาพประกอบจำลองเส้นเวลาประวัติศาสตร์
    \item สไลด์การสอนดิจิทัลและวิดีทัศน์สารคดีเกี่ยวกับนักวิทยาศาสตร์ผู้ค้นพบจุลินทรีย์ดิน
\end{itemize}

\textbf{การประเมินผล}
\begin{itemize}
    \item การประเมินจากการตอบคำถามและการมีส่วนร่วมในการอภิปรายในชั้นเรียน
    \item การตรวจแบบฝึกหัดท้ายบทและรายงานการวิเคราะห์โมเดลเศรษฐกิจชีวภาพ
\end{itemize}

\newpage

% ==========================================================
% เนื้อหาบทเรียน
% ==========================================================

\section{ความหมายและขอบข่ายของจุลชีววิทยาการเกษตร}
\index{จุลชีววิทยาการเกษตร (Agricultural Microbiology)}

\textbf{จุลชีววิทยาสำหรับการเกษตร (Agricultural Microbiology)} เป็นสาขาวิชาประยุกต์ของวิทยาศาสตร์ชีวภาพที่มุ่งเน้นศึกษาโครงสร้าง สรีรวิทยา พันธุศาสตร์ นิเวศวิทยา และปฏิสัมพันธ์ของจุลินทรีย์ชนิดต่างๆ ได้แก่ แบคทีเรีย (Bacteria) ฟังไจ (Fungi) แอกทิโนมัยซีท (Actinomycetes) อาร์เคีย (Archaea) สาหร่าย (Algae) และไวรัส (Viruses) ที่มีความเกี่ยวข้องกับระบบการผลิตทางการเกษตร ทั้งในมิติการส่งเสริมการเจริญเติบโตของพืช การฟื้นฟูดิน และการป้องกันกำจัดศัตรูพืช

\begin{definitionbox}{จุลชีววิทยาการเกษตร}
จุลชีววิทยาการเกษตร หมายถึง ศาสตร์ที่บูรณาการความรู้ด้านจุลชีววิทยากับเทคโนโลยีการผลิตพืชและสัตว์ เพื่อทำความเข้าใจปฏิสัมพันธ์ในระบบนิเวศดิน-พืช-จุลินทรีย์ และพัฒนาต่อยอดเป็นเทคโนโลยีชีวภาพที่สร้างผลผลิตทางการเกษตรอย่างยั่งยืนและเป็นมิตรต่อสิ่งแวดล้อม
\end{definitionbox}

\section{หมุดหมายสำคัญทางประวัติศาสตร์ของจุลชีววิทยาการเกษตร}
\index{ประวัติศาสตร์จุลชีววิทยา}

พัฒนาการของจุลชีววิทยาการเกษตรไม่ได้เกิดขึ้นอย่างโดดเดี่ยว หากแต่เชื่อมโยงกับวิวัฒนาการของกล้องจุลทรรศน์ เคมีเกษตร และนิเวศวิทยา โดยมีนักวิทยาศาสตร์ผู้บุกเบิกที่สำคัญดังนี้

ในคริสต์ศักราช 1676 \textbf{อันโทนี ฟาน เลเวนฮุก (Antonie van Leeuwenhoek, 1632--1723)} นักประดิษฐ์ชาวดัตช์ ได้ประดิษฐ์กล้องจุลทรรศน์เลนส์เดียวกำลังขยายสูง และเป็นมนุษย์คนแรกที่สังเกตพบจุลินทรีย์ในตัวอย่างน้ำและดิน ซึ่งเขาบันทึกเรียกว่า "สัตว์เซลล์เดียวขนาดเล็ก (Animalcules)"

ต่อมาในช่วงกลางคริสต์ศตวรรษที่ 19 \textbf{หลุยส์ ปาสเตอร์ (Louis Pasteur, 1822--1895)} นักเคมีและนักจุลชีววิทยาชาวฝรั่งเศส ได้พิสูจน์ให้เห็นว่า กระบวนการหมักไม่ได้เกิดขึ้นเองตามธรรมชาติ แต่เกิดจากกิจกรรมของสิ่งมีชีวิตขนาดเล็ก และปฏิเสธทฤษฎีการเกิดสิ่งมีชีวิตขึ้นมาเอง (Spontaneous Generation) อย่างสิ้นเชิง ซึ่งปูทางไปสู่วิธีการควบคุมจุลินทรีย์และการพาสเจอร์ไรซ์

ในด้านจุลชีววิทยาของดินโดยเฉพาะ นักวิทยาศาสตร์ผู้ได้รับการยกย่องให้เป็นบิดาแห่งจุลชีววิทยาดินคือ \textbf{เซอร์เกย์ วินโนกราดสกี (Sergei Winogradsky, 1856--1953)} นักจุลชีววิทยาชาวรัสเซีย ในปี ค.ศ. 1887 วินโนกราดสกีได้ค้นพบกระบวนการใช้พลังงานเคมีจากสารอนินทรีย์ (Chemoautotrophy) และค้นพบแบคทีเรียที่ทำหน้าที่ในกระบวนการไนตริฟิเคชัน (Nitrification) ได้แก่ \textit{Nitrosomonas} และ \textit{Nitrobacter} ตลอดจนแยกเชื้อแบคทีเรียตรึงไนโตรเจนแบบอิสระไม่ใช้ออกซิเจน \textit{Clostridium pasteurianum}

ในเวลาใกล้เคียงกัน ในปี ค.ศ. 1888 \textbf{มาร์ตินัส เบเจอรินค์ (Martinus Beijerinck, 1851--1931)} นักพฤกษศาสตร์และนักจุลชีววิทยาชาวดัตช์ ได้ประสบความสำเร็จในการแยกเชื้อแบคทีเรียตรึงไนโตรเจนที่อยู่ร่วมกับพืชตระกูลถั่วในปมราก คือ \textit{Rhizobium leguminosarum} และพัฒนาเทคนิคอาหารเลี้ยงเชื้อแบบคัดเลือก (Enrichment Culture Technique) ซึ่งเป็นรากฐานสำคัญในการแยกจุลินทรีย์จากธรรมชาติมาจนถึงปัจจุบัน

ในช่วงคริสต์ทศวรรษ 1940 \textbf{เซลแมน วาคส์แมน (Selman Waksman, 1888--1973)} นักจุลชีววิทยาดินชาวอเมริกันเชื้อสายยูเครน ได้ศึกษาแอกทิโนมัยซีทในดิน และค้นพบสารปฏิชีวนะสเตรปโตมัยซิน (Streptomycin) ซึ่งสามารถรักษาวัณโรคและโรคติดเชื้อแบคทีเรียได้ ส่งผลให้เขาได้รับรางวัลโนเบลสาขาสรีรวิทยาหรือการแพทย์ในปี ค.ศ. 1952 และกระตุ้นให้เกิดการค้นคว้าจุลินทรีย์ในดินเพื่อนำมาใช้ประโยชน์ทางการแพทย์และการเกษตรอย่างกว้างขวาง

\begin{figure}[H]
\centering
\begin{tikzpicture}[scale=0.95]
    % Main horizontal timeline line
    \draw[-{Stealth[scale=1.5]}, line width=2.5pt, color=agriGreen] (-0.3,0) -- (12.8,0);
    
    % Milestone 1: 1676 (Below)
    \draw[line width=1.5pt, color=agriGold] (0.8,0.25) -- (0.8,-0.25);
    \fill[agriLeaf] (0.8,0) circle (4pt);
    \node at (0.8, 0.55) [font=\footnotesize\bfseries\color{agriGreen}] {1676};
    \node at (0.8, -0.85) [align=center, font=\scriptsize\color{agriSlate}] {A. van Leeuwenhoek\\ค้นพบจุลินทรีย์ดิน\\(Animalcules)};
    
    % Milestone 2: 1860s (Above)
    \draw[line width=1.5pt, color=agriGold] (3.3,0.25) -- (3.3,-0.25);
    \fill[agriLeaf] (3.3,0) circle (4pt);
    \node at (3.3, -0.55) [font=\footnotesize\bfseries\color{agriGreen}] {1860s};
    \node at (3.3, 0.85) [align=center, font=\scriptsize\color{agriSlate}] {Louis Pasteur\\พิสูจน์กระบวนการหมัก\\ล้มล้างทฤษฎีเกิดขึ้นเอง};
    
    % Milestone 3: 1887 (Below)
    \draw[line width=1.5pt, color=agriGold] (6.0,0.25) -- (6.0,-0.25);
    \fill[agriLeaf] (6.0,0) circle (4pt);
    \node at (6.0, 0.55) [font=\footnotesize\bfseries\color{agriGreen}] {1887};
    \node at (6.0, -0.85) [align=center, font=\scriptsize\color{agriSlate}] {S. Winogradsky\\ค้นพบไนตริฟิเคชัน\\บิดาแห่งจุลชีววิทยาดิน};
    
    % Milestone 4: 1888 (Above)
    \draw[line width=1.5pt, color=agriGold] (8.7,0.25) -- (8.7,-0.25);
    \fill[agriLeaf] (8.7,0) circle (4pt);
    \node at (8.7, -0.55) [font=\footnotesize\bfseries\color{agriGreen}] {1888};
    \node at (8.7, 0.85) [align=center, font=\scriptsize\color{agriSlate}] {M. Beijerinck\\แยกเชื้อ Rhizobium\\Enrichment culture};
    
    % Milestone 5: 1940s-2026 (Below)
    \draw[line width=1.5pt, color=agriGold] (11.4,0.25) -- (11.4,-0.25);
    \fill[agriGold] (11.4,0) circle (5pt);
    \node at (11.4, 0.55) [font=\footnotesize\bfseries\color{agriGreen}] {1940s--2026};
    \node at (11.4, -0.85) [align=center, font=\scriptsize\color{agriSlate}] {S. Waksman $\to$ ยุค AI\\ปฏิชีวนะดิน $\to$ โอมิกส์\\เกษตรแม่นยำยั่งยืน};
\end{tikzpicture}
\caption{เส้นเวลาหมุดหมายสำคัญทางประวัติศาสตร์ของจุลชีววิทยาสำหรับการเกษตร}
\label{fig:intro_timeline}
\end{figure}

\section{จุลินทรีย์กับการขับเคลื่อนโมเดลเศรษฐกิจชีวภาพ}
\index{โมเดลเศรษฐกิจบีซีจี (BCG Model)}

ในบริบทของการพัฒนาที่ยั่งยืนแห่งสหัสวรรษ รัฐบาลไทยและประชาคมโลกได้ผลักดันโมเดลเศรษฐกิจสู่การพัฒนาที่ยั่งยืน (Bio-Circular-Green Economy: BCG Model) ซึ่งจุลินทรีย์การเกษตรมีบทบาทครอบคลุมทั้งสามเสาหลัก
\begin{enumerate}
    \item \textbf{เศรษฐกิจชีวภาพ (Bioeconomy)} มุ่งเน้นการใช้ประโยชน์จากทรัพยากรชีวภาพเพื่อสร้างมูลค่าเพิ่ม เช่น การใช้แบคทีเรียส่งเสริมการเจริญเติบโตของพืช (PGPR) เพื่อผลิตปุ๋ยชีวภาพทดแทนปุ๋ยเคมีราคาแพง
    \item \textbf{เศรษฐกิจหมุนเวียน (Circular Economy)} คำนึงถึงการนำวัสดุเหลือใช้ทางการเกษตรกลับมาใช้ประโยชน์ให้เกิดความคุ้มค่าสูงสุด ผ่านกระบวนการย่อยสลายทางชีวภาพ (Biodegradation) ของเชื้อราและแบคทีเรีย เช่น การทำปุ๋ยหมักจากทะลายปาล์มเปล่าและกากอ้อย
    \item \textbf{เศรษฐกิจสีเขียว (Green Economy)} มุ่งเน้นการพัฒนาที่รักษาสมดุลสิ่งแวดล้อม ลดการปล่อยก๊าซเรือนกระจก และรักษาความหลากหลายทางชีวภาพของดิน
\end{enumerate}

\begin{agribox}{จุลินทรีย์กับการเกษตรภาคตะวันออก}
จังหวัดจันทบุรีและกลุ่มจังหวัดภาคตะวันออกของประเทศไทย เป็นศูนย์กลางการผลิตผลไม้เมืองร้อนคุณภาพสูง เช่น ทุเรียน มังคุด และลำไย การประยุกต์ใช้จุลินทรีย์ปฏิปักษ์ เช่น เชื้อราไตรโคเดอร์มา (\textit{Trichoderma harzianum}) ในการควบคุมโรครากเน่าโคนเน่าจากเชื้อ \textit{Phytophthora palmivora} ถือเป็นตัวอย่างความสำเร็จของการเกษตรปลอดภัยที่ช่วยลดต้นทุนสารเคมีและเพิ่มความปลอดภัยแก่เกษตรกรและผู้บริโภค
\end{agribox}

\section*{คำถามทบทวนท้ายบทนำ}
\addcontentsline{toc}{section}{คำถามทบทวนท้ายบทนำ}
\begin{enumerate}
    \item จงอธิบายความสำคัญของเทคนิค Enrichment Culture ของ Martinus Beijerinck ต่อการพัฒนาจุลชีววิทยาการเกษตร
    \item เหตุใดการค้นพบกระบวนการ Chemoautotrophy ของ Sergei Winogradsky จึงเปลี่ยนมุมมองของนักวิทยาศาสตร์ต่อการหมุนเวียนธาตุในดิน
    \item จงวิเคราะห์ว่าการใช้จุลินทรีย์การเกษตรสามารถตอบโจทย์โมเดลเศรษฐกิจ BCG ในมิติเศรษฐกิจหมุนเวียน (Circular Economy) ได้อย่างไรบ้าง
\end{enumerate}

\clearpage
